\section{Self-Attention 的起点：Content Addressing 与 Parallel Reading}

上一章已经有 attention，却只用于 source--target interface。Self-attention 把 query、key、value 都来自同一序列，使每个 token 在一层内直接读取任意位置。它同时实现 global receptive field、position-parallel computation 和 multiplicative interaction；代价是 $n\times n$ score matrix 与顺序信息缺失。

\subsection{Attention-like Motif 早已存在于其他模态}

课程先展示 non-local means 与 texture synthesis，是为了反对“attention 完全凭空出现”的英雄叙事。\term{content-based addressing（内容寻址）}按当前 query 与候选内容的相似性决定读取位置，而不是按固定地址或 offset；它适合图像自相似、记忆检索和 source alignment。

\lecturefigure{slide-11-self-attention-other-modalities.jpg}{Non-local means 与 texture synthesis 中的内容寻址}{官方录像恢复幻灯片 V011；00:17:00}

\begin{knowledgebox}{读图：固定邻域与内容邻域是两种先验}
左侧 denoising 从整图寻找相似 patch，右侧 synthesis 从数据库寻找匹配纹理。先看被读取位置是否必须相邻，再看相似度由谁定义；这种机制和 attention 共享“query 决定连接”的结构，但传统算法的 feature/similarity 可能是手工设计，而 Transformer 端到端学习 projection。
\end{knowledgebox}

\subsection{Self-Attention：一句话中的 Token 互相更新}

上一小节说明了内容寻址为何能跨越固定邻域；这里进一步回答它怎样成为一层可训练的序列计算。读接下来的结构图时，先追踪单个 token 如何向全句发出查询，再观察所有 token 的查询为何可以批量并行完成。\term{self-attention（自注意力）}让序列中的每个位置产生 query、key、value，并用 query--key score 加权汇总 value。所有位置的 $Q,K,V$ 可由一次 matrix multiplication 生成，随后并行计算 score matrix。

\lecturefigure{slide-12-self-attention-structure.jpg}{Self-attention 让同一序列中的 token 直接交换信息}{官方录像恢复幻灯片 V012；00:17:27}

\begin{knowledgebox}{读图：一层内的图是 fully connected，但权重由数据决定}
中心词可以读取句中任意位置，不需经过多个 recurrent/convolutional step。先看每个输出是否都能连接所有输入，再看不同 query 得到不同权重；fully connected 指候选边存在，不代表所有边权相同，也不等于模型已知道语法关系。
\end{knowledgebox}

\lecturefigure{slide-13-self-attention-properties.jpg}{Self-attention 的三项设计属性}{官方录像恢复幻灯片 V013；00:17:36}

\begin{importantbox}{三项收益与一个隐含代价}
Self-attention 可并行、直接捕获 token--token interaction，并通过 dot product 提供显式 multiplicative gating。隐含代价是每对 token 都产生 score，标准实现的计算和 logit memory 随 $n^2$ 增长；“全连接”既是能力来源，也是 long-context 瓶颈。
\end{importantbox}

\subsection{原始目标还包括 Parallel Writing}

Transformer 的简化历史常说“RNN 太慢，所以换 attention”。课程补充了更激进的原始目标：既想并行读 source，也想并行写 target。\term{autoregressive decoding（自回归解码）}按 $p(y_t\mid y_{<t},x)$ 逐 token 缩小输出空间；若一次预测全部 token，则要同时学习哪些 token 可条件独立、先决定哪些 mode。

\lecturefigure{slide-14-transformer-original-motivation.jpg}{原始动机：并行读取，并尝试并行生成}{官方录像恢复幻灯片 V014；00:19:42}

自回归与完全非自回归 factorization 分别为：
\[
p_{\mathrm{AR}}(y\mid x)=\prod_{t=1}^{m}p(y_t\mid y_{<t},x),\qquad
p_{\mathrm{NAR}}(y\mid x)=\prod_{t=1}^{m}p(y_t\mid x).
\]
其中 $x$ 是输入，$y$ 是输出序列，$m$ 是输出长度。后一个公式假设给定 $x$ 后各 $y_t$ 条件独立，难以表示翻译中的多种合法措辞与彼此约束；实际 NAR 方法常加入 latent variable、iterative refinement 或分组顺序。

\teachervoice{讲者坦率地说 parallel decoding 的原始方案没有成功：自回归每次选词都会“弯曲”概率分布、逐步锁定 mode；非自回归模型还要学习决定顺序。团队最终保住了 parallel reading，却没有消灭 sequential generation。}

\begin{warningbox}{失败路线也是架构史的一部分}
Encoder parallelism 的成功不应被改写成最初唯一目标。记录 parallel writing 的失败，能解释为什么 decoder 仍需要 causal mask、为什么 speculative decoding 有现实价值，也提醒研究者区分理论并行与输出依赖结构。
\end{warningbox}

\subsection{本章小结}

Self-attention 把跨模态常见的 content retrieval 变成 sequence backbone，获得一层 global interaction 与并行表示学习。与此同时，fully parallel generation 因 mode ordering 和 conditional dependence 保持困难；Transformer 的胜利是一次筛选后的成功，而非所有原始目标同时实现。

